\verseline{}{我歌唱干戈与英雄：他从特洛伊海岸出发，\footnote{开篇 arma virumque cano 意为“我歌唱武器与英雄”。其中 virum 回应《奥德赛》首词 ἄνδρα（“那个人”）；arma 首先宣告《伊利亚特》式的战争题材； 《伊利亚特》首词 μῆνιν（“愤怒”）则从另一角度把战争置于史诗开篇。cano（“我歌唱”）回应荷马史诗开篇对缪斯的歌唱或讲述命令，维吉尔则先由诗人以第一人称宣布自己的歌唱。开篇兼具《伊利亚特》的战争题材与《奥德赛》的漂泊英雄；其中 arma 同时带出兵器与战争。primus（“首先”）强调埃涅阿斯开启拉丁—罗马谱系；第一卷稍后写到安忒诺耳先行抵达意大利并建立帕塔维乌姆。}}
\verseline{}{受命运驱逐，最初来到意大利、来到拉维尼乌姆的海岸；}
\verseline{}{他在陆地与海上久遭颠沛，}
\verseline{}{受天神之力驱迫，只因残酷的朱诺怀恨不忘；}
\verseline{5}{在战争中他也历经苦难，直至得以建起城邦，}
\verseline{}{将神祇带入拉提乌姆；由此诞生了拉丁民族、\footnote{此处 deos（神祇）主要指特洛伊的公共家神佩纳特斯。古代城邦可被理解为一个扩大的家族，也有自己的家神；埃涅阿斯的使命不仅是迁徙人口，也是把特洛伊的祭祀传统带入拉提乌姆。末三项依次指拉维尼乌姆、阿尔巴隆伽和罗马三个历史阶段。}}
\verseline{}{阿尔巴的先辈，以及高耸的罗马城墙。}
\verseline{}{缪斯啊，为我道出缘由：究竟哪一神意受了冒犯，\footnote{quo numine laeso 兼有两层含义：朱诺的神意受到阻挠，她的神性与尊严也受到冒犯。下文的 pietas 同时涵盖对神明、亲族、祖国和使命的责任。}}
\verseline{}{又为何悲愤，竟使那众神之后驱迫一位}
\verseline{10}{以虔敬著称的英雄，去历经如此多劫难，}
\verseline{}{去承受如此多劳苦？天神心中，竟有这样深的愤怒？}
\verseline{}{曾有一座古城，推罗的殖民者定居于此，}
\verseline{}{名为迦太基，隔海遥对意大利与台伯河口，}
\verseline{}{它财富丰饶，又尚武而极其剽悍；}
\verseline{15}{据说朱诺珍爱此地，胜过世间所有疆土，}
\verseline{}{就连萨摩斯岛也退居其后。她的兵器、\footnote{萨摩斯岛有著名的赫拉（罗马名朱诺）神庙。诗中朱诺却把迦太基置于萨摩斯之上，并在那里安置自己的兵器和战车；她对迦太基的偏爱也让罗马读者想到后来罗马与迦太基的长期敌对。}}
\verseline{}{她的战车皆在于此；若命运应允，}
\verseline{}{女神早已倾心谋划，要让它称霸万邦。}
\verseline{}{然而她曾听闻，特洛伊的血脉中将诞生一支后裔，}
\verseline{20}{终有一日会倾覆推罗的城堡；}
\verseline{}{一个统治四方、兵威骄盛的民族}
\verseline{}{将降临并毁灭利比亚：命运三女神正如此纺织。\footnote{拉丁文 Parcae 是罗马的命运女神，通常对应三位希腊摩伊赖：克洛托、拉刻西斯与阿特洛波斯。她们纺出命运的线；这里预言来自特洛伊血统的罗马人最终毁灭迦太基，暗指公元前 146 年第三次布匿战争的结局。}}
\verseline{}{萨图尔努斯之女对此心怀恐惧，又念及那场旧日之战——}
\verseline{}{她曾为挚爱的阿尔戈斯率先在特洛伊作战。}
\verseline{25}{不仅如此，那愤怒的缘由与狂暴的痛楚}
\verseline{}{尚未从她心头消散：帕里斯的判决，\footnote{帕里斯受命在朱诺、弥涅耳瓦与维纳斯之间裁定谁最美，最终把金苹果判给维纳斯，朱诺因此认为自己的美貌受辱。伽倪墨得斯则是被朱庇特带到奥林匹斯、取代朱诺之女赫柏担任斟酒者的特洛伊少年；他的荣耀又触动了朱诺对丈夫与特洛伊人的双重嫉恨。}}
\verseline{}{她因美貌受轻蔑而蒙受的屈辱，那为她憎恶的族裔，}
\verseline{}{以及被掳走的伽倪墨得斯所获的殊荣，依然深藏于心。}
\verseline{}{又因这些旧恨燃烧，她把那些在整个海上颠沛的}
\verseline{30}{特洛伊人——从达那俄斯人与无情的阿喀琉斯手中残存的人——}
\verseline{}{远远阻隔在拉提乌姆之外；许多年间，}
\verseline{}{他们受命运驱策，绕着所有海域漂流。}
\verseline{}{建立罗马民族，竟须承受如此巨大的艰难！}
\verseline{}{他们刚驶出西西里大地的视野，满怀喜悦地}
\verseline{35}{向深海扬帆，铜铸的船首破开咸腥的浪沫；}
\verseline{}{而朱诺，在胸中深藏着那永不愈合的创伤，}
\verseline{}{暗自思忖：“难道我这战败者，要就此放弃，}
\verseline{}{竟无法使特乌克罗斯人的国王远离意大利吗？\footnote{全诗以不同拉丁称谓指称同一批特洛伊人：Troes 为“特洛伊人”，Teucri 为“特乌克罗斯人”，Dardani ／ Dardanidae 为“达耳达诺斯人”，Phryges 为“佛里吉亚人”。称谓的变化取决于祖源、地域和说话立场。这里的 Teucrorum 指“特乌克罗斯人的国王”。}}
\verseline{}{是啊，命运阻挡了我。难道帕拉斯就能}
\verseline{40}{焚毁阿尔戈斯人的舰队，将他们淹没于深海，}
\verseline{}{只因俄琉斯之子埃阿斯一人的罪行与癫狂？}
\verseline{}{她亲自从云端掷下朱庇特迅疾的火焰，}
\verseline{}{又借狂风击碎战船、掀翻海面，}
\verseline{}{她用旋风卷起胸膛被刺穿、正喷吐着火焰的他，\footnote{小埃阿斯（俄琉斯之子）在特洛伊陷落时亵渎了帕拉斯的神庙，并凌辱或强掳女祭司卡珊德拉。帕拉斯因此用朱庇特的雷霆摧毁希腊舰队；朱诺拿这一先例质问，为何自己不能因一族之敌而惩罚整个特洛伊船队。}}
\verseline{45}{将他钉在尖锐的礁石上。}
\verseline{}{而我，堂堂众神之后，又是朱庇特的}
\verseline{}{姐姐兼妻子，竟同一个民族连续多年}
\verseline{}{征战不休！今后还有谁会敬奉朱诺的神威，}
\verseline{}{还有谁会前来祈求，在我的祭坛上敬献祭品？”}
\verseline{50}{女神在燃烧的心中反复盘算着这些，}
\verseline{}{来到了埃俄利亚——那风暴之乡，\footnote{埃俄利亚是风王埃俄卢斯的领域。埃俄卢斯把狂风关在山腹，以王权约束它们；稍后涅普顿平海的明喻又写一位有威望的人以言辞平息暴乱。两处都把自然力量写成需要统御的骚动。}}
\verseline{}{充满狂暴南风之地。在这里，国王埃俄卢斯}
\verseline{}{在巨大的洞穴中，以权威镇压着}
\verseline{}{相互角力的狂风与轰鸣的风暴，用锁链与囚牢将它们禁锢。}
\verseline{55}{它们愤懑不平，伴随着整座大山的巨大轰鸣，}
\verseline{}{围着栅门咆哮；而埃俄卢斯端坐于高高的要塞，}
\verseline{}{手握权杖，缓和它们的凶性，平息它们的怒火。}
\verseline{}{若非如此，那些狂风必定会席卷着大海、陆地}
\verseline{}{与深邃的天空，一同卷走，扫荡在长空之中。}
\verseline{60}{全能的父神正因担忧此事，才把它们藏入幽暗洞穴，}
\verseline{}{并在其上压下了一座座高耸的山脉；}
\verseline{}{他赐予它们一位国王，这位国王遵循着既定的法则，}
\verseline{}{懂得在受命之时收紧缰绳，或将其放松。}
\verseline{}{朱诺满怀恳求之意，向他吐露了这些话语：}
\verseline{65}{“埃俄卢斯啊，众神之父与人类之王既然赋予了你}
\verseline{}{平息波涛与乘风掀起巨浪的权力，}
\verseline{}{如今，一个与我为敌的民族正在第勒尼安海上航行，}
\verseline{}{他们要把伊利乌姆与战败的家神带往意大利：\footnote{Penates 指守护家庭储藏、炉火与延续的家神，范围比泛称的“众神”更具体。这里说“战败的家神”，把神祇也写成随特洛伊一起流亡的败者；他们同时又是未来罗马宗教连续性的凭证。}}
\verseline{}{将狂暴之力注入风中，倾覆并淹没他们的战船，}
\verseline{70}{或者将他们驱散，让躯体散落海面！}
\verseline{}{我有十四位身姿曼妙的宁芙，}
\verseline{}{其中最为美丽动人的得伊俄珀亚，}
\verseline{}{我将以恒久的婚姻使她与你结合，许她归你为妻，}
\verseline{}{为了酬谢你这般功勋，她将与你共度所有的岁月，}
\verseline{75}{并为你生下美丽的子嗣，让你成为父亲。”}
\verseline{}{埃俄卢斯如此答道：“王后，你想要什么，由你定夺；}
\verseline{}{遵命行事，则是我的本分。}
\verseline{}{是您为我赢得了我所得的这份王权，}
\verseline{}{为我赢得权杖与朱庇特的恩宠；是您准许我在众神宴席上落座，}
\verseline{80}{让我掌管云霭与风暴。”}
\verseline{}{话音刚落，他便倒转长矛，刺向空心山峰的}
\verseline{}{侧壁：狂风如同集结的阵列，}
\verseline{}{从敞开的缺口处倾泻而出，化作旋风席卷大地。}
\verseline{}{它们重重压向海面，从最深的海床处}
\verseline{85}{掀翻了整片汪洋。东风、南风与那裹挟着阵阵风暴的}
\verseline{}{西南风一齐呼啸，卷起狂涛巨浪拍向海岸。}
\verseline{}{紧接着响起众人的呼号与缆索的尖啸。}
\verseline{}{乌云瞬间夺走了天空与白昼，}
\verseline{}{特洛伊人的视线被遮蔽；漆黑的夜幕笼罩了海面。}
\verseline{90}{天穹轰鸣，高空闪烁着密集的电光，}
\verseline{}{万物都以迫在眼前的死亡威胁众人。}
\verseline{}{顷刻间，埃涅阿斯的四肢因冰冷的恐惧而瘫软；}
\verseline{}{他发出一声呻吟，向群星伸出双掌，}
\verseline{}{这样说道：“啊，三倍、四倍有福的人们，}
\verseline{95}{你们竟得以在父辈眼前、特洛伊高墙之下}
\verseline{}{迎面遇见死亡！哦，达那俄斯族中最勇猛的勇士，}
\verseline{}{提丢斯之子！我竟未能倒在伊利乌姆的原野，\footnote{“提丢斯之子”是狄俄墨得斯。本卷第 94—101 行中，埃涅阿斯在风暴里希望自己曾死在特洛伊战场；《奥德赛》5.306—312 中，遭遇海上风暴的奥德修斯也曾愿自己当初死在特洛伊。两处都让漂海英雄在迫近的海难中羡慕战场上的死者。}}
\verseline{}{未能让生命在你的右手之下流尽？}
\verseline{}{在阿喀琉斯的长矛下，凶猛的赫克托耳倒在那里，}
\verseline{100}{魁梧的萨尔佩冬也倒在那里；在那里，西摩伊斯河}
\verseline{}{曾在浪下卷挟多少战士的盾牌、头盔与强健的躯体！”}
\verseline{}{正当他喊出这番话时，风暴借着呼啸的北风}
\verseline{}{迎面击打船帆，把巨浪掀向群星。}
\verseline{}{船桨折断；紧接着船头偏转，将侧舷暴露给}
\verseline{105}{翻滚的波涛；一座陡峭的水之山轰然压下。}
\verseline{}{有人悬在浪峰；另一些人脚下巨浪裂开，}
\verseline{}{波涛间露出海底；狂潮裹着泥沙肆虐。}
\verseline{}{南风将三艘战船扭卷至隐藏的暗礁之上——}
\verseline{}{意大利人将这些横卧浪间的礁石称为“祭坛”——}
\verseline{110}{那是一道伏在水面下的巨大脊背；东风又将另外三艘}
\verseline{}{从深海卷入浅滩与流沙之中，景象惨不忍睹，}
\verseline{}{将它们砸向浅水，用沙堆将它们围困。}
\verseline{}{其中一艘载着吕喀亚人与忠诚的俄戎忒斯，}
\verseline{}{就在埃涅阿斯眼前，滔天巨浪从高处}
\verseline{115}{砸向船尾：舵手被震落，头朝下}
\verseline{}{翻滚入水；而海浪在原地将那艘船}
\verseline{}{旋转了三次，随后一个湍急的漩涡将其吞噬于海中。}
\verseline{}{浩瀚的旋涡中，零落的泅水者时隐时现，}
\verseline{}{战士的武器、木板与特洛伊的珍宝在波浪中漂浮。}
\verseline{120}{如今，伊利俄纽斯坚固的战船，勇士阿卡忒斯的战船，}
\verseline{}{以及载着阿巴斯和年迈阿勒忒斯的战船，}
\verseline{}{都被风暴压倒；它们侧壁的接缝松脱，}
\verseline{}{纷纷灌进凶猛的海水，从裂缝处开始崩解。}
\verseline{}{与此同时，涅普顿察觉到大海正被巨大的轰鸣声搅动，}
\verseline{125}{一场风暴被释放，宁静的水域从最深处}
\verseline{}{被翻腾起来，他心中极为震怒；他眺望深海，}
\verseline{}{把安详的面容探出水面。}
\verseline{}{他看到埃涅阿斯的舰队散落在整个海面上，}
\verseline{}{特洛伊人被波涛与天穹的崩塌所压迫，}
\verseline{130}{朱诺的愤怒与诡计没能瞒过她兄弟的双眼。}
\verseline{}{他将东风与西风召唤至面前，如此说道：}
\verseline{}{“你们对自己的出身竟有这般狂妄的自信？}
\verseline{}{如今，风啊，未经我的神圣旨意，你们竟敢搅乱}
\verseline{}{天空与陆地，掀起如此排山倒海的狂澜？}
\verseline{135}{我要把你们——但还是先平息这翻腾的波涛更为紧要。}
\verseline{}{此后，你们将用截然不同的惩罚偿还这次罪行。}
\verseline{}{加快你们逃离的步伐，去对你们的国王说：}
\verseline{}{海洋的统治权与这凶烈的三叉戟并不属于他，}
\verseline{}{而是经抽签归我。他占有的只是巨大的岩窟，}
\verseline{140}{东风啊，那是你们的居所；让埃俄卢斯在那座殿堂里}
\verseline{}{夸示自己的权威，统治那座紧闭的风之牢狱吧。”}
\verseline{}{话音未落，他比言语更迅速地平息了汹涌的波涛，}
\verseline{}{驱散了聚拢的乌云，带回了耀眼的阳光。}
\verseline{}{库摩托厄与特里同齐心协力，}
\verseline{145}{将战船从尖锐的礁石上推开；他亲自用三叉戟将船只撬起；}
\verseline{}{他打通了广阔的流沙，抚平了海面，}
\verseline{}{乘坐着轻盈的战车，滑行在波浪的顶端。}
\verseline{}{就像在一个庞大的人群中，常常会爆发}
\verseline{}{暴乱，卑贱的群氓在心中狂怒，}
\verseline{150}{火把与石块开始飞舞——狂热为他们送上武器；}
\verseline{}{此时，如果他们偶然望见一位因虔诚与功绩}
\verseline{}{而德高望重的人，他们便安静下来，侧耳倾听；\footnote{涅普顿平息海浪，被比作一位受人敬重的人以言辞制止暴乱。明喻中的 pietas 与开篇形容埃涅阿斯的 pius 同根。}}
\verseline{}{那人便用言语引导他们的心智，抚平他们的怒气，——}
\verseline{}{于是，大海的所有轰鸣也这般平息了，当这位父神}
\verseline{155}{眺望着水面，在开阔的苍穹下驾驭着战车，}
\verseline{}{转动马头，放开缰绳，乘车轻快飞驰。}
\verseline{}{疲惫不堪的埃涅阿斯众人，竭力驶向}
\verseline{}{最近的岸边，转向利比亚海岸。}
\verseline{}{在长长的海湾深处有一处隐秘的所在：一座岛屿的}
\verseline{160}{侧翼形成了一个天然的港口，从深海涌来的}
\verseline{}{每一道波浪都在这里被击碎，分流进入内收的港湾。}
\verseline{}{海港两侧巨岩耸立，两座峭峰}
\verseline{}{威临天际；在它们的峰脚下，}
\verseline{}{宽阔而安全的水域寂然无声；上方，闪亮的树林如布景般垂临，}
\verseline{165}{幽暗树丛覆下森然阴影。}
\verseline{}{正对面的崖壁下，有一处悬岩遮覆的洞穴，}
\verseline{}{里面有甘甜的泉水，也有天然岩石形成的座椅，}
\verseline{}{那是宁芙们的家园：在这里，疲惫的船只}
\verseline{}{不需要任何缆绳去固定，也不需要锚用弯曲的铁爪去咬合。}
\verseline{170}{埃涅阿斯带着从全队收拢的七艘船}
\verseline{}{驶入这里；怀着对陆地的深切渴望，}
\verseline{}{特洛伊人登上了期盼已久的沙滩，}
\verseline{}{把盐水淋漓的四肢放倒在海岸上。}
\verseline{}{阿卡忒斯最先用燧石敲击出火花，}
\verseline{175}{用树叶接住火种，在周围放上干燥的}
\verseline{}{引火物，让火苗在燃料中迅速燃起。}
\verseline{}{随后，被灾难拖得疲惫的人们取出海水泡坏的谷物}
\verseline{}{和磨谷的器具，将抢救回来的粮食}
\verseline{}{准备用火烘烤，并在石头上碾碎。}
\verseline{180}{与此同时，埃涅阿斯攀上岩峰，纵目搜寻}
\verseline{}{整片辽阔的海面，看能否望见安忒乌斯}
\verseline{}{在风中颠簸，以及佛里吉亚的双层桨船，}
\verseline{}{或卡皮斯，或高高船艉上卡伊库斯的盾牌。\footnote{“卡伊库斯的盾牌”原文为涵义较广的 arma，可泛指兵器、武装或战具。埃涅阿斯此时正从海上辨认失散船只；塞尔维乌斯（1.183）把它解释为系在船艉、用作识别标志的兵器，也列船具一解。译文取“盾牌”，是结合远距离辨认情境所作的具体化处理。}}
\verseline{}{视线中没有一艘战船，但他看见了三头雄鹿}
\verseline{185}{在海滩上徘徊；整个鹿群紧随其后，}
\verseline{}{排成长长的队伍在山谷中觅食。}
\verseline{}{他停下脚步，一把抓过忠诚的阿卡忒斯}
\verseline{}{所携带的武器——那张弓与疾飞的箭矢；}
\verseline{}{他先射倒三只鹿群的首领，它们高昂着那长着树枝般}
\verseline{190}{巨大鹿角的头颅，随后射向鹿群，用飞箭}
\verseline{}{把整个鹿群赶得在枝叶繁茂的树林间乱奔；}
\verseline{}{他没有停手，直到以胜者之势将七具巨大的}
\verseline{}{躯体放倒在地，使猎物的数量与他的战船相等。}
\verseline{}{接着他回到港口，将猎物分给所有的同伴。}
\verseline{195}{随后，他把好心的阿刻斯忒斯盛在酒坛里、}
\verseline{}{在特里纳克里亚海岸临别相赠的美酒，}
\verseline{}{分给众人，又用言语安抚悲痛的心：}
\verseline{}{“哦，同伴们——其实我们早已尝过苦难的滋味——}
\verseline{}{哦，经历了更沉重苦难的人们，神明也会让这些苦难终结。}
\verseline{200}{你们曾逼近斯库拉的狂暴与深处轰鸣的}
\verseline{}{礁石，也曾经历独眼巨人的岩壁：\footnote{斯库拉是海峡边吞噬水手的女怪；独眼巨人的岩壁指西西里附近波吕斐摩斯的居所。相关经历将在第三卷由埃涅阿斯倒叙。}}
\verseline{}{重振你们的勇气，放下悲哀的}
\verseline{}{恐惧：或许有朝一日，回忆起这些苦难也是一件乐事。}
\verseline{}{穿过种种遭际，穿过重重危境，}
\verseline{205}{我们正向着拉提乌姆进发；命运在那里向我们展示了}
\verseline{}{安宁的居所；在那里，特洛伊的王国将获准再次崛起。}
\verseline{}{坚持住，保全自己，以待顺境。”}
\verseline{}{他如此说道，尽管被巨大的忧虑折磨得心力交瘁，}
\verseline{}{他脸上装出希望的神色，却将深沉的悲痛压在心底。}
\verseline{210}{人们开始束紧衣袍处理猎物，准备即将到来的盛宴；}
\verseline{}{他们从肋骨上剥下兽皮，使血肉裸露；}
\verseline{}{有的人将肉切成碎块，把还在颤动的肉块串在烤肉叉上；}
\verseline{}{另一些人在海滩上架起铜锅，生起篝火。}
\verseline{}{随后，他们用食物唤回体力，散卧在草地上，}
\verseline{215}{尽情享用陈年的酒与肥美的鹿肉。}
\verseline{}{当宴席驱走饥饿，餐桌也被撤下，}
\verseline{}{他们在长谈中探寻着失去的同伴的下落，}
\verseline{}{他们在希望与恐惧之间迟疑：该相信同伴还活着，}
\verseline{}{还是已走到生命尽头，再也听不见呼唤。}
\verseline{220}{虔敬的埃涅阿斯在心中暗自悲叹，时而为勇猛的俄戎忒斯，}
\verseline{}{时而为阿米库斯，为吕库斯那残酷的命运，}
\verseline{}{以及勇敢的古阿斯与勇敢的克洛安图斯而哀伤。}
\verseline{}{当这一切告一段落，朱庇特从最高处的以太}
\verseline{}{俯视着帆影飞驰的海洋与平展的大地、}
\verseline{225}{海岸与分布广阔的万邦，他就这样伫立在天穹的巅峰，}
\verseline{}{将目光定格在利比亚的王国。}
\verseline{}{正当他在心中翻腾着这些思绪时，}
\verseline{}{维纳斯，面带愁容，闪烁的眼眸中充盈着泪水，}
\verseline{}{向他恳求道：“哦，以永恒的法则统御人类与众神之事，}
\verseline{230}{并以雷霆震慑万物的父神啊，}
\verseline{}{我的埃涅阿斯究竟对你犯下了何等重罪，}
\verseline{}{特洛伊人究竟做了什么，竟让他们在历经如此多死亡后，}
\verseline{}{仅仅为了意大利，整个世界竟对他们关闭？}
\verseline{}{你分明曾许诺，随着岁月流转，罗马人终将从这里诞生，}
\verseline{235}{从特乌克罗斯复苏的血脉中，将诞生统治者，}
\verseline{}{他们将以绝对的权威统治海洋与陆地，}
\verseline{}{父亲，究竟是什么改变了你的心意？}
\verseline{}{我曾以此慰藉特洛伊的覆亡与悲惨的废墟，}
\verseline{}{让相反的命运彼此抵偿；}
\verseline{240}{但如今，同样的厄运依旧纠缠着这些历经沧桑的人们。}
\verseline{}{伟大的神王啊，你将给他们的苦难定下怎样的终局？}
\verseline{}{安忒诺耳竟能从阿凯亚人的重围中逃脱，}
\verseline{}{深入伊利里亚的海湾，安全抵达利布尔尼亚人}
\verseline{}{最深处的国土，安然越过提马武斯河的源头，}
\verseline{245}{那里的水流伴随着山脉的巨大轰鸣，从九个喷口}
\verseline{}{倾泻而出，用咆哮的洪流淹没田野。\footnote{提马武斯河在亚得里亚海北岸，古人对其地下水道、泉口与潮汐现象有多种解释。九个出口与轰鸣洪流写出提马武斯河的奇景；安忒诺耳已经在意大利安居，埃涅阿斯却仍被拒于意大利之外。}}
\verseline{}{然而，他却在那里建起了帕塔维乌姆城，为特乌克罗斯人}
\verseline{}{找到了安居之所，赋予了民族新的名字，并悬挂起特洛伊的}
\verseline{}{武器；如今，他安居在宁静的和平中。}
\verseline{250}{而我们，你的后裔，你曾应允我们天穹的城堡，\footnote{“你曾应允我们天穹的城堡”暗示埃涅阿斯日后被接入神界。维纳斯称“我们”为朱庇特的后裔，是因为埃涅阿斯可经达耳达诺斯一系追溯至朱庇特；这一称呼依据神话谱系中的后裔关系。}}
\verseline{}{却（说来可悲！）失去了战船，仅仅因为某一个神明的愤怒}
\verseline{}{便遭到了背弃，被远远地隔离在意大利的海岸之外。}
\verseline{}{这难道是对虔诚的奖赏？你就是这样让我们重掌权杖的吗？”}
\verseline{}{人类与众神的缔造者对她微微一笑，}
\verseline{255}{带着那足以让天空与风暴放晴的容颜，}
\verseline{}{轻轻吻了女儿，随后如此说道：}
\verseline{}{“放下你的恐惧，库忒瑞亚：你子民的命运\footnote{“库忒瑞亚”是维纳斯的称号，得名于希腊南方的库忒拉岛；传说她从海中诞生后曾在附近登岸。这一古老称号维持了朱庇特神谕的庄严语体。}}
\verseline{}{为你稳固地保留着；你将看见这座城市与拉维尼乌姆}
\verseline{}{那应许的城墙，你将把伟大的埃涅阿斯}
\verseline{260}{高举至天穹的群星之间；我的心意未曾改变。\footnote{马克罗比乌斯（6.2.31）把《埃涅阿斯纪》开篇的风暴、维纳斯向朱庇特申诉，以及朱庇特以特洛伊人的未来安慰她视为一个整体，并说这一整组场景取自奈维乌斯《布匿战争》第一卷。奈维乌斯原诗今仅残存片段，这条古代见证仍很重要：维吉尔的开篇不只重写荷马，也主动承接早期拉丁历史史诗。}}
\verseline{}{他——既然这忧虑如此折磨着你，我便向你明言，}
\verseline{}{我将向更远处展开命运隐秘的书卷——}
\verseline{}{他将在意大利展开一场浩大的战争，}
\verseline{}{征服那些凶悍的民族，为人们建立制度与城墙，}
\verseline{265}{直到第三个夏天见证他在拉提乌姆的统治，}
\verseline{}{直到鲁图利人被征服后，三个驻营的冬天过去。\footnote{hiberna 可泛指冬季，但这里隐含 castra，即“冬营”。传统解释是：鲁图利人被征服后，埃涅阿斯又度过三个驻营的冬天，才建立拉维尼乌姆。}}
\verseline{}{而少年阿斯卡尼乌斯，如今又添上尤卢斯这一名号，}
\verseline{}{——当伊利乌姆国祚尚存时，他曾叫伊卢斯，——\footnote{阿斯卡尼乌斯又名伊卢斯、尤卢斯；尤卢斯这一名字使罗马尤利乌斯家族得以把祖先追溯到埃涅阿斯之子。}}
\verseline{}{他将以君王的身份，走完三十轮岁月，}
\verseline{270}{随后他会将王国从拉维尼乌姆的居所迁出，}
\verseline{}{用强大的力量筑起阿尔巴隆伽。\footnote{阿尔巴隆伽（Alba Longa）是传说中阿斯卡尼乌斯建立的城市，罗马王族谱系的重要中继点。原诗把地名拆开为 Longam... Albam；这里既是地名，也可直读为“绵长的阿尔巴”。}}
\verseline{}{在这里，赫克托耳的同族将整整统治三百个年头，}
\verseline{}{直到一位王族的女祭司——伊利娅，}
\verseline{}{因玛尔斯而有孕，生下一对双胞胎。}
\verseline{275}{此后，披着黄褐色狼皮、满怀喜悦的\footnote{nutricis（“乳母的”）指哺育罗慕路斯与雷穆斯的母狼。塞尔维乌斯（1.275） 将 tegmen 解作牧人披用的狼皮；罗慕路斯在这里被写成披着狼皮的建城者。}}
\verseline{}{罗慕路斯将接管这个民族，并建起玛尔斯的}
\verseline{}{城墙，用他自己的名字将他们命名为罗马人。}
\verseline{}{对他们，我不设疆界的限制，也不设时间的尽头；}
\verseline{}{我已赐予他们无尽的统治权。甚至那暴戾的朱诺，}
\verseline{280}{那个如今用恐惧折磨着海洋、陆地与天空的女神，}
\verseline{}{也将改变她的谋划，转向更好的道路；她将与我一同}
\verseline{}{庇佑罗马人——世界的主人、身披托加的民族。}
\verseline{}{这便是我的旨意。随着一个个五年祭期流逝，那个时代终将到来，\footnote{lustrum 原指监察官任期结束时举行的净化祭，通常为五年，因而也可泛指一段岁月。这里的“五年祭期”也带有朱庇特神谕中的罗马制度色彩。}}
\verseline{}{那时，阿萨剌科斯的家族将把佛提亚与辉煌的迈锡尼\footnote{阿萨剌科斯是特洛伊祖先，埃涅阿斯为其曾孙。佛提亚、迈锡尼与阿尔戈斯分别关联阿喀琉斯、阿伽门农与狄俄墨得斯；预言说昔日战败的特洛伊后裔将反过来统治这些希腊地区。}}
\verseline{285}{压制为奴，并在战败的阿尔戈斯称王。}
\verseline{}{一位出身高贵的特洛伊人——恺撒——将要诞生，\footnote{通说多指奥古斯都；亦有尤利乌斯·恺撒或有意兼指二人的解释。后文的东方战利品、受人祈祷及神化更切合奥古斯都在阿克提乌姆战役（即亚克兴战役）后的形象，而“恺撒—尤利乌斯—尤卢斯”的连续命名又让两位恺撒与特洛伊谱系重叠。}}
\verseline{}{他的帝国将以海洋为界，他的荣耀将直达星辰，——}
\verseline{}{尤利乌斯，这正是传承自伟大的尤卢斯的名字。}
\verseline{}{有朝一日，你将在天国安心地迎接他，满载着东方的战利品；}
\verseline{290}{人们也将在祈祷中呼唤他。}
\verseline{}{那时，暴戾的时代将随着战争的平息而变得温和；}
\verseline{}{古老的信义女神与维斯塔，还有奎里努斯与兄弟雷穆斯，\footnote{奎里努斯是罗慕路斯神化后的名字。神谕让他与在建城冲突中死去的兄弟雷穆斯共同颁布法律；这对在建城传说中反目的兄弟在预言里重新并列。“古老”在拉丁文中同时修饰信义女神与维斯塔，唤起理想化的古代德性。}}
\verseline{}{将共同颁布律法；战争那可怕的大门\footnote{“战争之门”通常指雅努斯神庙之门：战争时开启，和平时关闭。奥古斯都曾以关闭此门宣示和平；门内被锁住的“狂乱”尤其暗指罗马内战。}}
\verseline{}{将被坚固的铁栓封锁；而那渎神的狂乱，}
\verseline{295}{将坐在凶器之上，被锁在门内，}
\verseline{}{双手被一百个铜结绑在背后，用沾满鲜血的嘴发出可怖的咆哮。”}
\verseline{}{他如此说着，并从高天派下了迈亚之子，\footnote{“迈亚之子”是墨丘利。他是朱庇特的信使，也主管道路、商贸和往来；朱庇特派他先使迦太基人平息敌意，为特洛伊人进入城市创造条件。}}
\verseline{}{为了让这片土地与新建的迦太基城堡向特乌克罗斯人}
\verseline{}{敞开好客的大门，以免对命运一无所知的狄多}
\verseline{300}{将他们驱逐出境：他在浩瀚的空中用翅膀划过，}
\verseline{}{迅速降落在利比亚的海岸上。}
\verseline{}{他立即执行了命令，在神明的意志下，}
\verseline{}{布匿人放下了凶猛的心性；尤其是王后，}
\verseline{}{以平静的意愿和仁慈的心接纳特乌克罗斯人。}
\verseline{305}{而虔敬的埃涅阿斯，在整夜翻腾着无数的思绪后，}
\verseline{}{当温煦的晨光初现，便决定走出去，}
\verseline{}{探查这片陌生土地，看看他被风吹到了哪一片海岸，}
\verseline{}{查明这片看似荒芜的土地上，住的究竟是人还是野兽，}
\verseline{}{并将探明的一切带回给同伴们。}
\verseline{310}{他将舰队隐藏在林间空地深处的岩洞下，}
\verseline{}{四周被树木和森然的阴影封住；}
\verseline{}{他只带阿卡忒斯一人作伴，向前走去，}
\verseline{}{手中挥动着两支宽刃长矛。}
\verseline{}{在森林深处，他的母亲迎面向他走来，}
\verseline{315}{她以少女的面容与装束出现，带着斯巴达少女的}
\verseline{}{武器，又如同色雷斯的哈耳帕吕刻，\footnote{哈耳帕吕刻是色雷斯的女猎手与战士，以奔跑迅疾著称。维纳斯把自己装扮成斯巴达少女般的女猎手，又被比作奔跑迅疾的色雷斯哈耳帕吕刻；她虽作凡人装束，声音与容貌仍让埃涅阿斯怀疑她是女神。}}
\verseline{}{将烈马驱使至疲惫，并以超越赫布鲁斯河的速度飞奔。}
\verseline{}{像女猎手一样，她在肩上习惯性地挂着一张轻巧的弓，}
\verseline{}{任由长发在风中铺展，}
\verseline{320}{膝盖裸露着，流动的衣襟用一个结收拢。}
\verseline{}{她先开口道：“喂，年轻人们，告诉我，}
\verseline{}{你们是否偶然看见我的某个姐妹在这里游荡，}
\verseline{}{她带着箭袋，披着斑驳的猞猁皮，}
\verseline{}{或者正高声呼喊着，追捕一头口吐白沫的野猪？”}
\verseline{325}{维纳斯如此说道；而维纳斯之子如此回答：}
\verseline{}{“我没有听到、也没有看到你的任何一位姐妹——}
\verseline{}{哦⋯⋯我该怎样称呼你，姑娘？因为你的面容}
\verseline{}{绝非凡人，你的声音也不似人类：哦，你必定是位神明——}
\verseline{}{你是福玻斯的姐妹？还是宁芙血脉中的一员？——}
\verseline{330}{无论你是谁，愿你开恩，减轻我们的劳苦，}
\verseline{}{告诉我们，我们究竟身在哪片天空下，}
\verseline{}{被抛到大地的哪一方海岸。我们对这里的人民与地域一无所知，}
\verseline{}{在风与巨浪驱赶下茫然来到此地；}
\verseline{}{我必亲手让许多祭牲倒在你的祭坛前。”}
\verseline{335}{维纳斯回答道：“我绝不敢领受这般殊荣；}
\verseline{}{推罗少女的习俗便是背负箭袋，}
\verseline{}{并用紫色的高帮靴紧紧绑住小腿。}
\verseline{}{你眼前是布匿人的国土、推罗人和阿革诺耳的城市；\footnote{阿革诺耳是腓尼基王族的祖先，常被视为腓尼克斯之父；“阿革诺耳的城市”因此就是以祖先之名称呼推罗—腓尼基城市。维吉尔交替使用“推罗人”“腓尼基人”“布匿人”等名称，所指都是这批腓尼基移民。}}
\verseline{}{但这片领土属于利比亚人——一个在战争中难以驯服的民族。}
\verseline{340}{狄多统治着这个帝国，她从推罗城出逃，}
\verseline{}{逃避她的哥哥。这场冤屈说来话长，}
\verseline{}{曲折也说来话长；但我只讲事情的梗概。}
\verseline{}{她的丈夫是叙凯乌斯，拥有腓尼基最丰厚的土地，}
\verseline{}{也是这不幸女子深深爱过的人；}
\verseline{345}{她的父亲在初婚的吉兆中将纯洁的她许配给他。}
\verseline{}{但推罗的王权掌握在她哥哥}
\verseline{}{皮格马利翁的手中；他作恶凶残，胜过所有人。}
\verseline{}{狂乱横亘在二人之间。皮格马利翁}
\verseline{}{亵渎神明，在祭坛前，又被黄金之爱蒙蔽，}
\verseline{350}{暗中用刀剑杀死毫无防备的叙凯乌斯，}
\verseline{}{全不顾妹妹的情意；他把这桩罪行隐瞒许久，}
\verseline{}{编造种种谎言，以空洞的希望欺骗为爱情憔悴的妹妹。}
\verseline{}{但在睡梦中，她未葬丈夫的幻影亲自显现，}
\verseline{}{以奇异的方式抬起惨白的面庞，}
\verseline{355}{显出那残酷的祭坛与被刀剑刺穿的胸膛，}
\verseline{}{揭开家中全部隐秘的罪行。}
\verseline{}{他催促她加快逃跑的步伐，离开祖国，}
\verseline{}{并从地下揭示了古老的宝藏以资助她的旅途，}
\verseline{}{那是无人知晓的无数白银与黄金。}
\verseline{360}{被这一切震惊的狄多，开始筹划逃亡，召集同伴：}
\verseline{}{聚来的是那些痛恨暴君}
\verseline{}{或深怀恐惧的人；他们夺取了碰巧备妥的船只，}
\verseline{}{装满了黄金：贪婪的}
\verseline{}{皮格马利翁的财富就这样被运上海去；而领头的是一名女子。}
\verseline{365}{他们来到了你如今将看到的这些地方，那里巨大的}
\verseline{}{城墙与新迦太基正在崛起的城堡清晰可见，}
\verseline{}{他们买下一块土地；那地方由此得名柏尔萨，}
\verseline{}{其大小刚好是一张牛皮所能圈住的范围。\footnote{柏尔萨建城传说称，狄多获准购买“一张牛皮所能围住”的土地，便把牛皮裁成极细长条，围出足以筑城的大片区域。Byrsa 又与希腊语“牛皮”相近；这个双关故事产生于希腊语世界的地名解释传统。}}
\verseline{}{但你们到底是谁？你们从哪个海岸而来，}
\verseline{370}{又要前往何方？”面对她的询问，}
\verseline{}{他叹息着，从胸膛深处吐出这样的话语：}
\verseline{}{“哦，女神，如果我从最初的源头开始讲述，}
\verseline{}{而你有闲暇听完我们劳苦的编年史，}
\verseline{}{那么，在我讲完之前，晚星早已关闭奥林匹斯，收起白昼。}
\verseline{375}{我们来自古老的特洛伊，如果特洛伊的名字}
\verseline{}{曾偶然飘入过你的耳中，在穿越了各种海洋之后，}
\verseline{}{风暴偶然地将我们推向了利比亚的海岸。}
\verseline{}{我是虔敬的埃涅阿斯，我带着从敌人手中抢救出的家神，}
\verseline{}{随我的舰队一同航行，我的名声传到了苍天之上。}
\verseline{380}{我寻找祖国意大利；我的族系源自至高的朱庇特。\footnote{Conte 2019 将 et genus ab Ioue summo（“我的族系源自至高的朱庇特”）置于方括号；该版校勘记认为这几个词可能由第六卷第123行补入。意大利被称为“祖国”，与达耳达诺斯的意大利祖源传统相联。}}
\verseline{}{我曾率二十艘战船驶入佛里吉亚海域，}
\verseline{}{在神明母亲的指引下，追随命运所示的道路；}
\verseline{}{如今，遭波涛与东风重创，只剩七艘。}
\verseline{}{而我在异乡无人知晓、贫乏无援，流浪在利比亚的荒野，}
\verseline{385}{被逐出欧罗巴与亚细亚。”维纳斯没有让他继续抱怨下去，}
\verseline{}{而是在他悲痛的倾诉中插话道：}
\verseline{}{“无论你是谁，既然你来到了推罗的城市，}
\verseline{}{我相信，你呼吸着生命的空气，并未遭天神憎恶。}
\verseline{}{只管继续前行，把你的脚步迈向女王的门槛，}
\verseline{390}{因为我向你宣告：你的同伴已归来，你的舰队已被寻回，}
\verseline{}{在北风的转向下，它们已被送入安全的水域，}
\verseline{}{除非父母传给我的占卜之术全然无用。}
\verseline{}{看那排成阵列的十二只欢跃的天鹅，\footnote{十二只天鹅对应风暴中失散的十二艘船；鹰是朱庇特的鸟，天鹅则与维纳斯关系密切。天鹅被鹰驱散后重新结队、部分落地而部分即将降落，预示船队有的已经入港，有的正在驶入。}}
\verseline{}{朱庇特的神鹰曾从以太的领域滑翔而下，}
\verseline{395}{在开阔的天空中驱散了它们；如今排成长长的行列，}
\verseline{}{它们似乎正在降落，或正在俯瞰已被它们占据的土地：}
\verseline{}{正如它们归来时，用呼啸的翅膀嬉戏着，}
\verseline{}{在天空中盘旋成群，并发出欢快的鸣叫，}
\verseline{}{载着你那些年轻同伴的船只也正是这样：}
\verseline{400}{有的已停在港内，有的正满帆驶入港口。}
\verseline{}{只管前行，沿着道路所引的方向迈步。”}
\verseline{}{她说完转过身去，玫瑰色的颈项闪耀光辉，}
\verseline{}{不死的长发从头顶吐出神圣的}
\verseline{}{芬芳；长袍一直垂落到脚底，}
\verseline{405}{步态间显出真正的女神。当他}
\verseline{}{认出那是自己的母亲时，他对着她远去的背影高声呼喊：}
\verseline{}{“为何连你也要如此残忍，屡屡用虚假的幻影}
\verseline{}{嘲弄你的儿子？为何不让我们右手相握，}
\verseline{}{去聆听和回答彼此真实的声音？”}
\verseline{410}{他如此责备着，将脚步迈向城墙：}
\verseline{}{但维纳斯用幽暗的气流将这两个前行的人包裹，}
\verseline{}{女神在他们周围覆上一层厚厚的云雾，如同披风，}
\verseline{}{以免有任何人看见他们，或触碰他们，}
\verseline{}{以免有人制造延误，或是盘问他们此行的缘由。}
\verseline{415}{她自己则高飞着离去，回到了帕福斯，}
\verseline{}{满怀喜悦地重返自己的居所；那里有她的神庙，\footnote{帕福斯在塞浦路斯，是维纳斯的重要崇拜中心。萨巴香指来自阿拉伯南部萨巴地区的珍贵乳香。百座祭坛燃着乳香，新鲜花环散发香气；维纳斯已经从凡人伪装回到自己的神殿。}}
\verseline{}{一百座祭坛燃着萨巴香，透着温热，吐出新鲜花环的芬芳。}
\verseline{}{与此同时，他们沿着小路指引的方向，加快了步伐。}
\verseline{}{很快他们登上了那座巨大的山丘，}
\verseline{420}{它高高地悬在城市上方，从高处俯视着对面的城堡。}
\verseline{}{埃涅阿斯惊叹于这座庞大的城市——此地从前只有茅屋，}
\verseline{}{他惊叹于那些城门、喧闹的人群与铺砌的街道。}
\verseline{}{推罗人正奋力劳作：有人延展城墙，}
\verseline{}{建造城堡，用双手滚上巨石；}
\verseline{425}{有人为房屋挑选地基，用犁沟圈定范围。}
\verseline{}{他们选定法律、官员与神圣的元老院；\footnote{Conte 2019 将第426行置于方括号；该版校勘记列出删去本行、移至第368行之后，以及保留本行的意见。此行把法律、官员与元老院写入初建的迦太基。}}
\verseline{}{在这里，有的人在挖掘港口；在那里，另一些人}
\verseline{}{在为剧院铺设深深的地基，并从岩壁上}
\verseline{}{凿下巨大的圆柱，作为未来舞台高耸的装饰。}
\verseline{430}{就像在初夏阳光照耀下的鲜花盛开的原野上，}
\verseline{}{劳作驱使着蜜蜂，当它们带领着家族中成年的}
\verseline{}{后代外出，或是当它们将流动的蜂蜜}
\verseline{}{压紧，用甜美的花蜜将蜂房胀满，}
\verseline{}{或是接过归来者的重担，或是排成阵列，}
\verseline{435}{将那些懒惰的雄蜂驱赶出蜂巢：}
\verseline{}{工作在沸腾，芬芳的蜂蜜散发着百里香的气息。}
\verseline{}{“哦，何其幸运的人们，他们的城墙已经在升起！”}
\verseline{}{埃涅阿斯感慨道，仰视着城市的屋顶。}
\verseline{}{他裹在云雾中进入人群——说来令人惊叹——}
\verseline{440}{穿行其中，混入众人，却没有谁看见他。}
\verseline{}{在城市的中央有一片树林，阴凉而宜人，}
\verseline{}{正是在这里，那些被波涛和旋风抛掷的布匿人，}
\verseline{}{在尊贵的朱诺所指引的地点，挖出了}
\verseline{}{一个征兆——一颗烈马的头；因为这意味着这个民族\footnote{烈马之头是迦太基的吉兆，象征尚武与富足，马的图像也常见于迦太基钱币。后来埃涅阿斯在意大利见到白母猪；两处征兆分别同迦太基和特洛伊人在意大利的建城前景相连。}}
\verseline{445}{将在战争中卓尔不群，并在世世代代中享有丰饶的物产。}
\verseline{}{在这里，西顿的狄多正在为朱诺建造一座宏伟的神庙，}
\verseline{}{里面摆满了丰厚的献礼，充满了女神的神力，}
\verseline{}{青铜门槛在台阶上高高升起，横梁由青铜联结，}
\verseline{}{青铜大门在铰链上发出沉重的鸣响。}
\verseline{450}{在这片树林中，眼前这陌生的景象第一次}
\verseline{}{平息了他的恐惧，在这里，埃涅阿斯第一次}
\verseline{}{敢于期望得救，并对自己饱受打击的境遇重拾信心。}
\verseline{}{因为，当他在巨大的神庙中细细环顾四周，}
\verseline{}{等待着王后，当他惊叹于这座城市的气象，}
\verseline{455}{惊叹于工匠们竞相施展的手艺与营造的辛劳时，}
\verseline{}{他看到了伊利乌姆的战役被按顺序描绘，}
\verseline{}{那是早已被名声传遍了整个世界的战争，}
\verseline{}{他看到了阿特柔斯之子们，看到了普里阿摩斯，以及对双方都冷酷无情的阿喀琉斯。}
\verseline{}{他停下脚步，流着泪说道：“阿卡忒斯啊，如今在这世上}
\verseline{460}{还有哪个地方，哪片土地不曾装满我们的苦难？}
\verseline{}{看呐，普里阿摩斯！即便在这里，荣耀也有它应得的回报；}
\verseline{}{这里也有世事的眼泪；人间种种，触动人心。\footnote{sunt lacrimae rerum 的关键在 rerum：它既可作宾语性属格，意为“有为世事而流的泪”，也可感受为主语性或归属性关系，即“世事自身有泪”。mortalia 则是开放的中性复数，可指凡人之事、人间遭际，也可在语境中偏向苦难。埃涅阿斯把壁画看作迦太基人能够同情特洛伊苦难的证据，因而相信自己会得到庇护。紧接着的 pictura inanis 又显出这种安慰的限度：图像就在眼前，所画的往事与死者却不可复得。}}
\verseline{}{放下你的恐惧吧；这传扬的名声将为你带来某种庇护。”}
\verseline{}{他如此说着，以那空幻的画面滋养心神，}
\verseline{465}{一再叹息，任泪水成流，浸湿面庞。}
\verseline{}{因为他看到，当他们在帕加马周围鏖战时，}
\verseline{}{希腊人在这里溃退，特洛伊的青年紧追不舍，}
\verseline{}{而在另一边，佛里吉亚人正在逃窜，盔缨高耸的阿喀琉斯驾着战车逼近。}
\verseline{}{就在不远处，他流着泪认出瑞索斯那}
\verseline{470}{覆着雪白篷布的帐幕；营中人在初睡时便遭出卖，\footnote{色雷斯王瑞索斯得到预言：若他的白马吃到特洛伊草料、饮过克珊托斯河水，特洛伊便不会陷落。狄俄墨得斯与乌利塞斯趁他初次入睡时杀死他并夺走马匹；这一夜袭故事见《伊利亚特》第十卷。}}
\verseline{}{提丢斯之子在那里大肆杀戮，浑身染血，}
\verseline{}{把那些烈马赶回营中，在它们尚未}
\verseline{}{尝到特洛伊的草料、饮过克珊托斯河水之前。}
\verseline{}{画面的另一处，特洛伊卢斯失去武器，正在逃跑，}
\verseline{475}{不幸的男孩，与阿喀琉斯交锋却力不相敌，}
\verseline{}{他被战马拖拽着，仰面挂在空荡荡的战车上，\footnote{特洛伊卢斯是普里阿摩斯的幼子，主要属于《伊利亚特》之外的特洛伊战争传说。壁画把他写成不敌阿喀琉斯的少年，被失控的战马拖曳；他仍握着缰绳，倒转的长矛一路划过地面。}}
\verseline{}{手中却依然紧攥着缰绳；他的脖子和头发}
\verseline{}{被拖在地上，倒转的矛尖在尘土中划出痕迹。}
\verseline{}{与此同时，伊利乌姆的妇女走向不肯相助的帕拉斯神庙，}
\verseline{480}{她们披散着头发，带着供奉的长袍，}
\verseline{}{充满悲哀地恳求着，用手掌捶打着胸膛；}
\verseline{}{但女神转过身去，双眼始终低垂于地面。}
\verseline{}{阿喀琉斯将赫克托耳的尸体绕着伊利乌姆的城墙拖拽了三圈，}
\verseline{}{并以黄金为价，出售这具失去生命的躯体。\footnote{《伊利亚特》中阿喀琉斯把赫克托耳的尸体绕帕特罗克洛斯坟墓拖行，维吉尔却写成绕特洛伊城三圈。普里阿摩斯后来以丰厚赎礼换回儿子的尸体；原诗用“出售”这一冷酷动词，把赎尸写得近似交易。}}
\verseline{485}{看到这一幕，他从胸膛的最深处发出一声悲号，}
\verseline{}{当他看到那战利品，看到那战车，看到朋友的尸体，}
\verseline{}{以及手无寸铁的普里阿摩斯伸出的双掌。}
\verseline{}{他也认出了与阿凯亚诸将混战的自己，}
\verseline{}{东方的军阵，还有黑肤门农的兵甲。}
\verseline{490}{狂怒的彭忒西勒亚率领着亚马逊人的军队，}
\verseline{}{带着新月形的盾牌，在万军之中如火焰般燃烧，\footnote{彭忒西勒亚是亚马逊女王，在《伊利亚特》正文结束后的传说中率军援救特洛伊，最终死于阿喀琉斯之手。新月形盾牌与单侧裸胸是古代艺术中亚马逊战士的常见标识。}}
\verseline{}{在袒露的乳房下，系着金色的束带，}
\verseline{}{这位女战士敢与男子交锋——她还是个少女。}
\verseline{}{当这些画面让达耳达诺斯的埃涅阿斯感到震惊时，}
\verseline{495}{当他惊愕不已，凝神注视、久久不移时，}
\verseline{}{王后，容貌绝美的狄多，向着神庙}
\verseline{}{走来，一大群青年紧紧簇拥在她的周围。}
\verseline{}{就像在欧罗塔斯河岸，或库恩图斯山脊上，}
\verseline{}{狄安娜引领舞队，一千名山宁芙追随她，\footnote{欧罗塔斯河流经斯巴达，库恩图斯山位于提洛岛，二者都与狄安娜／阿耳忒弥斯崇拜相连。这里狄多被比作高出千名宁芙的狄安娜；《奥德赛》6.102—110 也以阿耳忒弥斯比拟即将接待奥德修斯的瑙西卡娅。}}
\verseline{500}{从两旁聚集；她肩负箭袋，}
\verseline{}{行进时高过所有女神；}
\verseline{}{喜悦悄然漫过拉托娜的胸膛。}
\verseline{}{狄多也是如此，她满怀喜悦地走在人群之中，}
\verseline{}{督促工程，营建未来的王国。}
\verseline{505}{随后，在女神的殿门前，在神庙中央的穹顶下，}
\verseline{}{她由武装侍卫环绕，高高地坐在王座上。}
\verseline{}{她为人们裁断案件、颁布法律，将工程的劳作}
\verseline{}{公平地分配，或是通过抽签来决定：}
\verseline{}{就在这时，埃涅阿斯忽然看见在大群人簇拥下走来的\footnote{安忒乌斯等特洛伊人前来求见狄多，迦太基人聚拢在他们周围。}}
\verseline{510}{安忒乌斯、塞耳革斯图斯与勇敢的克洛安图斯，}
\verseline{}{以及其他特乌克罗斯人——他们曾被海上那黑暗的旋风}
\verseline{}{吹散，并被远远地卷向了其他的海岸。}
\verseline{}{他与阿卡忒斯同时惊呆，被狂喜}
\verseline{}{与恐惧击中；他们热切地渴望与同伴右手相握；}
\verseline{515}{但未知的情势扰乱了他们的心智。}
\verseline{}{他们深藏在云幕中，继续暗中观察，}
\verseline{}{想看看这些人遭遇了什么命运，把舰队留在了哪个海岸，}
\verseline{}{又为何而来；因为这些从各艘船上被选出的人，}
\verseline{}{正大声呼唤着祈求恩典，走向神庙。}
\verseline{520}{他们进去，获准当面陈词后，}
\verseline{}{最年长的伊利俄纽斯怀着平静的心情，如此开口道：}
\verseline{}{“哦，王后，朱庇特赐予了你建立新城，}
\verseline{}{并以正义约束骄傲的民族，}
\verseline{}{我们这些可悲的特洛伊人，被狂风卷过所有的海洋，}
\verseline{525}{祈求你，让那可怕的火焰远离我们的船只，}
\verseline{}{怜惜一个虔诚的民族，并仔细审视我们的遭遇。}
\verseline{}{我们来此，不是要用刀剑劫掠利比亚人的家园，}
\verseline{}{也不是要把夺来的战利品赶向海岸；}
\verseline{}{我们心中没有这种意图，战败者也没有这般傲慢。}
\verseline{530}{有一个地方，希腊人称之为赫斯珀里亚，\footnote{赫斯珀里亚意为“西方之地”，希腊人常以此称意大利。俄诺特里亚人是意大利南部的古老居民；诗中又用传说领袖伊塔卢斯解释“意大利”之名，这是古代常见的名称起源说。}}
\verseline{}{那是一片古老的土地，武力强盛而泥土肥沃；}
\verseline{}{俄诺特里亚人曾居住于此；如今有传言说，他们的后裔}
\verseline{}{用他们领袖的名字，将这个民族称为意大利。}
\verseline{}{那本是我们的航向：}
\verseline{535}{当带来风暴的俄里翁突然从浪涛上升起，\footnote{俄里翁既是猎人神话人物，也是星座；其升落在古代航海经验中与风暴季节相联。“带来风暴的俄里翁升起”以猎人星座的人格形象写海上恶候。}}
\verseline{}{将我们推入隐秘的浅滩，又在肆虐的南风中，}
\verseline{}{在海浪的翻滚与不可逾越的礁石间，将我们远远地}
\verseline{}{吹散；我们只有少数人游到了你的海岸。}
\verseline{}{这究竟是什么样的族群？什么祖国竟容许}
\verseline{540}{如此野蛮的习俗？我们连在海滩容身都被拒绝；}
\verseline{}{他们挑起战争，不许我们踏上最初的陆地。}
\verseline{}{即使你们轻视人类与凡人的武器，}
\verseline{}{至少也该敬畏那些铭记神律所容与所禁的神明。}
\verseline{}{我们的国王是埃涅阿斯，在正义上，}
\verseline{545}{在虔诚上，在战争与武器上，无人能出其右。}
\verseline{}{如果命运保全了这个人，如果他仍在呼吸着以太的}
\verseline{}{空气，尚未倒在残酷的阴影中，}
\verseline{}{我们便无所畏惧；你若抢先施恩，}
\verseline{}{也绝不会后悔。在西西里的地区，同样有着城邦}
\verseline{550}{与武装力量，还有源自特洛伊血脉、声名卓著的阿刻斯忒斯。}
\verseline{}{请允许我们把被风暴重创的舰队拖上岸，}
\verseline{}{允许我们在森林中准备木材并削制船桨：}
\verseline{}{如果命运允许我们找回同伴和国王，并继续驶向}
\verseline{}{意大利，那么我们将满怀喜悦地向着意大利和拉提乌姆进发；}
\verseline{555}{但如果这拯救的希望已被夺走，如果你——特乌克罗斯人最好的父老——}
\verseline{}{已被利比亚的海洋吞没，若尤卢斯也再无生还的希望，}
\verseline{}{那么，至少让我们去寻找西坎尼亚的海峡与那已准备好的安居之所，}
\verseline{}{我们就是从那里被吹来的，让我们去投奔阿刻斯忒斯国王吧。”}
\verseline{}{伊利俄纽斯如此说道；与此同时，所有达耳达诺斯人}
\verseline{560}{都齐声响应。}
\verseline{}{接着，狄多微微低垂着脸庞，简短地回答道：}
\verseline{}{“特乌克罗斯人，从你们的心中解开恐惧，收起你们的忧虑吧。}
\verseline{}{残酷的时局与王国的新建迫使我}
\verseline{}{采取这样的手段，用卫兵在四处守卫我的边疆。}
\verseline{565}{谁会不知道埃涅阿斯一族，谁会不知道特洛伊城，}
\verseline{}{不知道他们的勇武与英雄，或那场大战的烈火？}
\verseline{}{我们布匿人的心胸还没有迟钝到如此地步，}
\verseline{}{太阳驾车巡天，也没有远远绕开推罗城。}
\verseline{}{无论你们是向往伟大的赫斯珀里亚与萨图尔努斯的田野，}
\verseline{570}{还是选择厄律克斯的边界与阿刻斯忒斯国王，}
\verseline{}{我都将在护卫下让你们安全离去，并用财富资助你们。}
\verseline{}{或者，你们愿意与我平等地在这片王国里定居？}
\verseline{}{我正在建立的这座城市就是你们的；把你们的船只拖上岸吧；}
\verseline{}{无论特洛伊人还是推罗人，我都一视同仁。}
\verseline{575}{啊，但愿埃涅阿斯国王本人，也被同样的南风驱赶，}
\verseline{}{来到了这里！我一定会派出可靠的人沿着海岸}
\verseline{}{去搜索，我将命令他们巡视利比亚的尽头，}
\verseline{}{看看他是否在遭遇海难后，正在哪片森林或城市中流浪。”}
\verseline{}{这些话语让勇敢的阿卡忒斯与父亲埃涅阿斯}
\verseline{580}{精神振奋，他们早已迫不及待地想要冲破云层的}
\verseline{}{束缚。阿卡忒斯率先对埃涅阿斯说道：}
\verseline{}{“女神之子，你的心中此刻升起了怎样的念头？}
\verseline{}{你看到一切都很安全，舰队和同伴都已找回。}
\verseline{}{只有一人下落不明，我们曾亲眼看着他在海浪中}
\verseline{585}{沉没；其余的一切都正如你母亲所言。”}
\verseline{}{他刚说完这些话，周围包裹着他们的云雾便突然}
\verseline{}{裂开，消散在清澈的以太之中。}
\verseline{}{埃涅阿斯站在那里，在明亮的光芒中闪耀，}
\verseline{}{他的面庞与宽阔的肩膀宛如神明；因为他的母亲亲自}
\verseline{590}{将鬈发的秀美、青春的玫瑰光泽}
\verseline{}{以及眼中明朗的神采吹给了儿子：}
\verseline{}{正如巧手为象牙增添光彩，又如白银或}
\verseline{}{帕罗斯大理石被嵌入黄金。}
\verseline{}{随后他这样对着王后说道，突然出现在所有人的}
\verseline{595}{意料之外：“你要找的人就在你面前，}
\verseline{}{我是特洛伊的埃涅阿斯，从利比亚的海浪中死里逃生。}
\verseline{}{哦，唯有你怜悯特洛伊那难以言说的苦难，}
\verseline{}{唯有你，将你的城市与家园，与我们这些从达那俄斯人}
\verseline{}{手中残存下来，早已在陆地与海洋的各种灾厄中耗尽心力、}
\verseline{600}{一无所有的人们共享；给你相称的报答，}
\verseline{}{并不在我们能力的范围之内，狄多，也不在任何}
\verseline{}{散落在这广阔世界各处的达耳达诺斯人的能力之中。}
\verseline{}{如果还有任何神明看顾着虔敬的人们，如果正义在任何地方}
\verseline{}{仍有分量，也还有自知正直的心，愿神明们\footnote{Conte 2019 在第604行采用主格 iustitia；该版校勘记将 si quid iustitia est 释为“如果正义还有任何效力／分量”。此处问正义是否仍有力量。}}
\verseline{605}{赐予你相配的奖赏。是怎样喜悦的时代孕育了你？}
\verseline{}{是怎样伟大的父母生下了如此的女儿？}
\verseline{}{只要河流依然奔向大海，只要阴影仍在山坡间}
\verseline{}{巡行，只要天穹仍在哺育着群星，}
\verseline{}{你的荣耀、你的名字与你的颂声就将永远流传，}
\verseline{610}{无论哪片土地召我前往。”他说完，伸出右手}
\verseline{}{握住朋友伊利俄纽斯，左手又握住塞雷斯图斯，}
\verseline{}{随后是其他人，勇敢的古阿斯与勇敢的克洛安图斯。}
\verseline{}{西顿的狄多，首先是对眼前这景象感到震惊，}
\verseline{}{随后又惊叹于这男人如此巨大的遭际。她开口说道：}
\verseline{615}{“女神之子，是怎样的灾厄让你历经如此巨大的危险？}
\verseline{}{是怎样的力量将你推向这荒蛮的海岸？}
\verseline{}{难道你就是那个埃涅阿斯，那个滋养万物的维纳斯}
\verseline{}{在佛里吉亚的西摩伊斯河畔，为达耳达诺斯的安喀塞斯生下的儿子？}
\verseline{}{我确实记得，特乌克罗斯曾来到西顿，\footnote{这里的 Teucer 是人物特乌克罗斯：他是忒拉蒙与特洛伊公主赫西俄涅之子，站在希腊一方参战，战后被父亲逐出故土，后来在塞浦路斯建立新的萨拉弥斯；他的母系血统使他仍可自称出自特洛伊。}}
\verseline{620}{他被赶出了父辈的疆界，在贝卢斯的帮助下}
\verseline{}{寻找着新的王国；那时，我的父亲贝卢斯正在洗劫富庶的}
\verseline{}{塞浦路斯，并以胜利者的姿态将其置于统治之下。}
\verseline{}{从那时起，我就已知晓特洛伊城的陨落，}
\verseline{}{知晓你的名字，以及佩拉斯吉的诸王。}
\verseline{625}{他本人虽曾是敌人，却盛赞特乌克罗斯人，}
\verseline{}{还说自己也出自特乌克罗斯人古老的血脉。}
\verseline{}{所以来吧，年轻人们，进入我们的屋檐下吧。}
\verseline{}{相似的命运也曾使我饱经劳苦，}
\verseline{}{最后才让我在这片土地上安顿下来。}
\verseline{630}{我并非不懂苦难；我学会了援助不幸的人。”}
\verseline{}{她如此说道；同时引领着埃涅阿斯进入那皇家}
\verseline{}{殿堂，同时下令在神庙中献上祭品以示尊崇。}
\verseline{}{与此同时，她也并未忘记派人前往海岸，送给他的同伴们}
\verseline{}{二十头公牛，一百头背生粗鬃的}
\verseline{635}{大猪，一百只肥美的羔羊与它们的母亲，}
\verseline{}{作为今日的礼物与欢乐。}
\verseline{}{宫殿内室闪耀着王家的奢华，}
\verseline{}{已经布置妥当；他们在殿堂中央准备盛宴：}
\verseline{}{那里有工艺精湛、闪耀着华贵紫色的织锦，}
\verseline{640}{桌上摆着巨大的银器，黄金上雕刻着}
\verseline{}{先祖们英勇的壮举，那是一系列从民族}
\verseline{}{古老起源开始、经无数英雄延续而下的漫长历史。}
\verseline{}{埃涅阿斯（因为父亲对儿子的爱让他的心无法平静）}
\verseline{}{派阿卡忒斯迅速返回船上，}
\verseline{645}{让他把这些事告诉阿斯卡尼乌斯，并把他带到城墙内；}
\verseline{}{这位慈父所有的关切都集中在阿斯卡尼乌斯身上。}
\verseline{}{此外，他还下令带来那些从伊利乌姆的废墟中抢救出来的}
\verseline{}{礼物：一件用黄金与图案刺绣得挺括的披风，}
\verseline{}{一条边缘织有藏红花色茛苕图案的面纱，}
\verseline{650}{这些曾是阿尔戈斯的海伦的饰物，\footnote{海伦把这些华美织物从迈锡尼带到特洛伊，伴随的是她与帕里斯“不被允许的婚姻”。埃涅阿斯此刻把引发特洛伊灾难的婚礼遗物赠给狄多，而这些礼物又将点燃另一段致命爱情，物件本身带有强烈的情节反讽。}}
\verseline{}{当她前往帕加马，去寻求那不被允许的婚姻时，从迈锡尼}
\verseline{}{带出的，这是她母亲勒达送给她的一件奇妙的礼物；}
\verseline{}{此外，还有普里阿摩斯最年长的女儿伊利俄涅\footnote{伊利俄涅是普里阿摩斯最年长的女儿，传统中嫁给色雷斯王波吕墨斯托耳。她的权杖、珍珠项链与双层王冠把特洛伊王室的残余尊荣带入狄多宫廷。}}
\verseline{}{曾经握过的权杖，一条}
\verseline{655}{挂满珍珠的项链，以及一顶镶嵌着宝石与黄金的双层王冠。}
\verseline{}{阿卡忒斯带着这些吩咐，迅速向着船只走去。}
\verseline{}{但库忒瑞亚却在心中盘算着新的计谋，新的策略：}
\verseline{}{让丘比特改变他的容貌与脸庞，}
\verseline{}{去顶替甜蜜的阿斯卡尼乌斯，用这些礼物}
\verseline{660}{点燃王后，使她发狂，将爱的火焰缠绕进她的骨髓之中；}
\verseline{}{因为她不信任这难测的王室，也不信任口是心非的推罗人；}
\verseline{}{凶狠的朱诺使她焦灼，忧虑总在夜幕降临时卷土重来。}
\verseline{}{因此，她用这些话语对着长着翅膀的爱神说道：}
\verseline{}{“孩子，你是我唯一的力量，我强大的权柄，}
\verseline{665}{孩子，你连至高父神那劈碎提丰的雷霆都不放在眼里，\footnote{提丰（Typhoeus）是反抗奥林匹斯诸神的巨怪，被朱庇特以雷霆击败。维纳斯说丘比特连这雷霆也不畏惧，以此强调他的力量。}}
\verseline{}{我向你寻求庇护，恳求你的神力相助。}
\verseline{}{你的哥哥埃涅阿斯，因为严酷的朱诺对他心怀仇恨，}
\verseline{}{在海上被抛向四方海岸，}
\verseline{}{这你是知道的，你也常因我的痛苦而悲伤。}
\verseline{670}{如今，腓尼基女子狄多留住了他，用温柔的话语}
\verseline{}{延缓他的行程；我担心朱诺的这场款待}
\verseline{}{将走向何方；在这样关键的转折点，她绝不会停手。}
\verseline{}{因此，我打算用计谋先发制人，用爱的火焰}
\verseline{}{包围这位王后，免得她受哪位神明的力量左右而变心，}
\verseline{675}{而是要让她和我一样，被对埃涅阿斯深深的爱所束缚。}
\verseline{}{你现在听听我的想法，看看如何能做到这一点。}
\verseline{}{这位王族的男孩、我最深的牵挂，正准备}
\verseline{}{应他亲爱父亲的召唤，带着那些从海洋与特洛伊的烈火中}
\verseline{}{留存下来的礼物，前往西顿的城市：}
\verseline{680}{我将把沉睡的他藏在高高的库忒拉}
\verseline{}{或是伊达利乌姆的神圣居所中，}
\verseline{}{使他无从察觉这场计谋，也不能中途介入。}
\verseline{}{你要用计谋伪装成他的容貌，只需这一个夜晚，}
\verseline{}{以一个男孩的身份，换上那个男孩为人熟知的脸庞，}
\verseline{685}{这样，当充满喜悦的狄多把你抱入怀中，}
\verseline{}{在王家的餐席与吕埃乌斯的酒浆间，\footnote{“吕埃乌斯”是巴克科斯的称号，常解释为“解忧者”或“松解者”，此处借称葡萄酒。维纳斯让丘比特趁狄多在宴饮、拥抱与亲吻中放松戒备时，把隐秘的火注入她体内。}}
\verseline{}{当她给你拥抱，并印下甜蜜的亲吻时，}
\verseline{}{你便能将那隐秘的火焰吹入，用爱情的毒药迷惑她。”}
\verseline{}{爱神遵从了他亲爱母亲的吩咐，卸下了}
\verseline{690}{翅膀，满心欢喜地迈出了尤卢斯的步伐。}
\verseline{}{而维纳斯则将平静的睡意灌遍阿斯卡尼乌斯的四肢，}
\verseline{}{女神将他抱在怀里，带到了伊达利乌姆高高的}
\verseline{}{树林中，那里柔软的马郁兰用花朵与甜美的}
\verseline{}{阴凉环绕着他，散发着芬芳。}
\verseline{695}{此时丘比特正听从着命令，带着皇家的礼物}
\verseline{}{前往推罗人那里，在阿卡忒斯的引领下显得十分快活。}
\verseline{}{当他抵达时，王后已经在华丽的帷幔下}
\verseline{}{在金色的长榻上安顿好自己，并坐在了中央。}
\verseline{}{随后父亲埃涅阿斯与特洛伊的青年们}
\verseline{700}{相继到来，斜倚在铺开的紫色织锦上。}
\verseline{}{侍从为他们的双手倾水，从篮中端出面包，}
\verseline{}{并递上剪去绒毛的光滑餐巾。}
\verseline{}{内殿里有五十名侍女，负责依次备下}
\verseline{}{接连不断的菜肴，并在火焰中礼敬家神；}
\verseline{705}{还有一百名其他侍女，以及同等数量、年龄相仿的男仆，}
\verseline{}{负责在桌上摆满菜肴，并摆上酒杯。}
\verseline{}{同样，推罗人也从各处欢乐的门口涌入，}
\verseline{}{聚集在这里，受邀斜倚在饰纹的长榻上。}
\verseline{}{他们惊叹于埃涅阿斯的礼物，惊叹于尤卢斯，}
\verseline{710}{惊叹于神明那闪耀的面容与他那伪装的话语，}
\verseline{}{惊叹于那件披风与绣着藏红花色茛苕的面纱。}
\verseline{}{尤其是那个不幸的女人，那注定要遭遇灾厄的}
\verseline{}{腓尼基女子，她的心灵无法得到满足，她在凝视中燃烧，}
\verseline{}{被这个男孩和那些礼物同样地打动着。}
\verseline{715}{他依偎过埃涅阿斯，挂在他的颈上，}
\verseline{}{满足了那位受蒙蔽之父深沉的慈爱，}
\verseline{}{便走向王后。她的目光、她的整个心胸}
\verseline{}{都依恋着他，时而把他搂在怀中；狄多浑然不知，}
\verseline{}{坐在她膝上的竟是何等强大的神。}
\verseline{720}{然而他牢记着阿喀达利亚母亲的嘱托，\footnote{“阿喀达利亚”是维纳斯的罕见称号，传统上联系她与阿喀达利亚泉；塞尔维乌斯（1.720）还把它与“刺痛、忧虑”的词源联想相连。丘比特奉命以对埃涅阿斯的“活的爱”压过亡夫叙凯乌斯的记忆。}}
\verseline{}{开始渐渐抹去她心中叙凯乌斯的记忆，并试图用鲜活的爱情}
\verseline{}{先占住她久已沉寂的心神、早已不惯于爱的心。}
\verseline{}{宴席初歇，餐桌撤下之后，}
\verseline{}{人们摆出巨大的调酒缸，为酒浆戴上花冠。\footnote{vina coronant 字面为“给酒加冠”，酒可代指盛酒的器皿。塞尔维乌斯（1.724） 分列斟至器沿、以花环装饰酒器两解；按后一解，花冠装饰的是盛酒器。}}
\verseline{725}{谈笑声在高屋下升起，沿宽阔的}
\verseline{}{厅堂回荡；点燃的灯盏从金色的藻井上}
\verseline{}{垂下，火把用火焰战胜了黑夜。}
\verseline{}{这时，王后要来一只缀满宝石与黄金的沉重祭酒盘，}
\verseline{}{将纯酒斟满其中；贝卢斯以及贝卢斯的所有后裔}
\verseline{730}{向来都使用这只酒盘；随后大厅里安静下来：}
\verseline{}{“朱庇特，因为人们说你为宾主立下了律法，}
\verseline{}{愿你让这一天，成为推罗人与从特洛伊远道而来的人们共同的欢乐，}
\verseline{}{愿我们的子孙后代铭记这一天。}
\verseline{}{愿赐予欢乐的巴克科斯与慈佑的朱诺同在；}
\verseline{735}{还有你们，哦，推罗人，请满怀善意地庆祝这场聚会吧。”}
\verseline{}{她说完，便将一些酒洒在桌上作为祭酒，\footnote{patera 是宽而浅的祭酒盘，可倾洒少量酒液，也可直接饮用。狄多先向朱庇特献酒，因为朱庇特保护宾主之间的神圣义务；她又在祷词中把推罗人与特洛伊人的后代合称“我们的后人”，预示她对永久联合的愿望。}}
\verseline{}{献酒之后，她先以双唇轻触酒杯，}
\verseline{}{然后催促着递给比提亚斯；他毫不迟疑地饮尽}
\verseline{}{杯中起沫的酒浆，从满盛的金杯里豪饮而下，}
\verseline{740}{随后其他贵族也纷纷效仿。留着长发的伊俄帕斯}
\verseline{}{弹奏起金色的里拉琴；伟大的阿特拉斯曾教导他。\footnote{伊俄帕斯是狄多宴上的歌者。阿特拉斯既是托举天穹的巨神，也在古代传统中被视为星象知识的教师；伊俄帕斯的歌涉及日月运行、生命起源、天气与星座。}}
\verseline{}{他歌唱游移的月亮与太阳的亏蚀；}
\verseline{}{歌唱着人类与兽群的起源；歌唱着雨水与火焰的由来；}
\verseline{}{歌唱大角星、多雨的许阿得斯星群与大小二熊；\footnote{大角星位于牧夫座；许阿得斯位于金牛座，古人认为它的升起预示降雨；“大小二熊”即大熊座与小熊座，拉丁文亦将它们拟为牛牵引的车辆。Conte 2019 仍将第 744 行正常排入正文；该版校勘记指出此行与第三卷第 516 行完全相同，并记录 Ribbeck 的删疑意见。接下来的诗句（1.745—746）则追问冬季昼短夜长的自然成因。}}
\verseline{745}{歌唱着冬日的太阳为何如此匆忙地要去大洋中沐浴，}
\verseline{}{又是什么阻滞，拖住了迟迟不去的长夜。}
\verseline{}{推罗人掌声再起，特洛伊人也跟着喝彩。}
\verseline{}{狄多也以变换不尽的谈话拉长黑夜，}
\verseline{}{那不幸的人，正饮下绵延不尽的爱，}
\verseline{750}{不断询问普里阿摩斯，反复询问赫克托耳；}
\verseline{}{时而问奥罗拉之子带着怎样的武器前来，\footnote{奥罗拉之子是埃塞俄比亚王门农，他在特洛伊战争后期援救特洛伊，并与阿喀琉斯决战。狄多先追问普里阿摩斯、赫克托耳、门农、狄俄墨得斯与阿喀琉斯的片段，最后要求埃涅阿斯从木马之计开始连续讲述。埃涅阿斯随后在异邦女王的宫廷中叙述战争与漂泊；《奥德赛》9—12 中奥德修斯也曾在费阿刻斯人的宫廷作长篇自述。}}
\verseline{}{时而问狄俄墨得斯的战马如何，时而问阿喀琉斯何等强大。}
\verseline{}{“来吧，客人，从最初的源头为我们讲起，”}
\verseline{}{她说道，“讲讲达那俄斯人的诡计，你人民的灾厄，}
\verseline{755}{以及你的流浪；因为如今已是第七个夏天，}
\verseline{}{带着你在所有的陆地和波涛间漂泊。”}
